\subsection{Separation of convex sets in a locally convex space}

PROPOSITION 4. --- Let A be a closed non-empty convex set in a locally convex space E and let K be a compact non-empty convex set in E, that does not meet A. Then there exists a closed hyperplane H that strictly separates A from K.

For there exists an open convex neighbourhood V of 0 in E such that A + V and K + V do not meet (GT, II, § 4.3, prop. 4). As A + V and K + V are convex and open in E, prop. 1 of II, p.~37 shows that there exists a closed hyperplane H that strictly separates A + V from K + V, and a fortiori A from K.

Remark. --- In a Hausdorff locally convex space E, let A and B be two closed non-empty convex sets that are disjoint, if E is finite dimensional then there exists a closed hyperplane that separates A from (II, p.~78, exerc. 13); but this conclusion is not necessarily true when E is of infinite dimension (II, p.~78, exerc. 10 and 11).

COROLLARY 1. --- In a locally convex space, every closed convex set A is the intersection of the closed half-spaces which contain it.

In fact, for every point \(x \notin \mathbf{A}\) , there exists a closed hyperplane that separates \(x\) strictly from \(\mathbf{A}\) (using prop. 4).

COROLLARY 2. --- In a Hausdorff locally convex space, every compact convex set A is the intersection of the closed half-spaces which contain it and which are determined by support hyperplanes of A.

For, let \(x_{0} \notin A\) ; \(\{x_{0}\}\) is closed, therefore there exists a closed hyperplane H which separates \(x_{0}\) strictly from A (prop. 4); let \(f(x) = \alpha\) be an equation of H (f a continuous linear form) and suppose that \(f(x) > \alpha\) for all \(x \in A\) . If we put \(\gamma = \inf_{x \in A} f(x)\) , the half-space defined by \(f(x) \geqslant \gamma\) contains A, is determined by the support hyperplane of equation \(f(x) = \gamma\) , and does not contain \(x_{0}\) ; whence the corollary.

It is possible that a closed convex set that is not compact and has no interior point, in a locally convex space, does not have any closed support hyperplane (II, p.~86, exerc. 18 : cf.~also V, p.~71, exerc. 11).

COROLLARY 3. --- In a locally convex space, the closure of each linear variety M is the intersection of the closed hyperplanes that contain M.

For all \(x \notin \overline{\mathbf{M}}\) , let \(\mathbf{H}\) be a closed hyperplane that separates \(x\) strictly from \(\overline{\mathbf{M}}\) ;

thus \(\overline{M}\) is parallel to H; the closed hyperplane \(H_{1}\) , containing \(\overline{M}\) and parallel to H does not contain x. The corollary follows.

COROLLARY 4. --- Let C be a closed convex set in a locally convex space E. A subset A of E is contained in C, if, and only if, for every real valued continuous affine function u in E such that \(u(x) \geqslant 0\) for all x in C, we have \(u(y) \geqslant 0\) for all y in A.

The condition is obviously necessary. Conversely we show that it is sufficient; if a point \(x \in A\) is not contained in C, there exists a closed hyperplane of equation \(f(z) = \alpha\) separating x strictly from C; if we suppose for example that \(f(x) < \alpha\) , then the continuous affine function \(u = f - \alpha\) contradicts the hypotheses.

COROLLARY 5. --- In a locally convex space E, the closure of each convex cone C of vertex 0 is the intersection of closed half-spaces containing C determined by closed hyperplanes that pass through 0.

For \(\overline{\mathbf{C}}\) is a convex cone of vertex 0 (II, p.~13, prop. 14). For \(x \notin \overline{\mathbf{C}}\) , there exists a closed hyperplane H that separates \(x\) strictly from \(\overline{\mathbf{C}}\) (prop. 4). It is now just necessary to apply the following lemma:

Lemma 1. --- If a cone A, with vertex 0, is contained in an open half-space determined by a hyperplane H, then it is contained in a closed half-space determined by a hyperplane \(H_{0}\) , that is parallel to H and passes through 0.

Let \(f(z) = \alpha\) with \(\alpha < 0\) be an equation of H, so that \(f(z) = 0\) is the equation of \(H_0\) . If there exists \(z \in A\) such that \(f(z) < 0\) , then there would exist \(\lambda > 0\) , such that \(f(\lambda z) = \alpha\) , and as \(\lambda z \in A\) , this would contradict the hypothesis.

